% Lean module: MonadFormal.Information.Likelihood
\subsection{Likelihoods under input coarse-graining}
\label{sec:input-likelihood}

Let $X,M,Y$ be measurable spaces, $q:X\to M$ measurable,
$\mu\in\mathcal P(X)$, and $K,L:M\rightsquigarrow Y$ Markov
channels. Put $\lambda=q_*\mu$ and $\widehat q(x,y)=(q(x),y)$.
The two channels define pairs of joint laws
\[
 \rho_X=\mu\otimes(K\circ q),\quad
 \sigma_X=\mu\otimes(L\circ q),\qquad
 \rho_M=\lambda\otimes K,\quad
 \sigma_M=\lambda\otimes L.
\]
Here the subscripts identify the input carrier, not a marginal.

\begin{lemma}[Pullback of a joint density]
\label{lem:input-density}
The joint laws satisfy $\widehat q_*\rho_X=\rho_M$ and
$\widehat q_*\sigma_X=\sigma_M$. If a measurable
$f:M\times Y\to[0,+\infty]$ satisfies $\rho_M=f\sigma_M$, then
\[
 \rho_X=(f\circ\widehat q)\sigma_X.
\]
\end{lemma}
\begin{proof}
For a measurable rectangle $A\times B\subseteq M\times Y$,
\[
 (\widehat q_*\rho_X)(A\times B)
 =\int_{q^{-1}(A)}K(q(x),B)\,d\mu(x)
 =\int_A K(m,B)\,d\lambda(m).
\]
This proves the first pushforward identity; replace $K$ by $L$ for
the second. Hence $(f\circ\widehat q)\sigma_X$ has the same finite
total mass as $f\sigma_M=\rho_M$.

Fix a measurable $B\subseteq Y$. The density identity, tested on
every measurable input set $A$, gives
\[
 \int_A\int_B f(m,y)\,L(m,dy)\,d\lambda(m)
 =\int_A K(m,B)\,d\lambda(m).
\]
Uniqueness of nonnegative densities therefore gives
\[
 \int_B f(m,y)\,L(m,dy)=K(m,B)
 \quad\text{for $\lambda$-almost every $m$}.
\]
Pull this equality back through $q$ and integrate over an arbitrary
measurable subset of $X$. The resulting equality on rectangles proves
the claimed identity between finite measures. The exceptional set may
depend on $B$; no simultaneous conditional-density identity is assumed.
\leanproof{MonadFormal.Information.map_input_joint,MonadFormal.Information.withDensity_input_joint}
\end{proof}

\begin{theorem}[Relative entropy under input coarse-graining]
\label{thm:input-relative-entropy}
On arbitrary measurable carriers,
\[
 D_{\mathrm{KL}}(\rho_X\Vert\sigma_X)
 =D_{\mathrm{KL}}(\rho_M\Vert\sigma_M).
\]
\end{theorem}
\begin{proof}
Data processing under $\widehat q$ gives the inequality from right
to left. If $\rho_M\not\ll\sigma_M$, its relative entropy is
$+\infty$, so that inequality already forces equality.

Otherwise take $f=d\rho_M/d\sigma_M$.
By \Cref{lem:input-density}, $\rho_X\ll\sigma_X$ with density
$f\circ\widehat q$. For $z\ge0$, put
$\phi(z)=z\log z+1-z$, with $0\log0=0$.
The nonnegative-integral formula for relative entropy yields
\[
 \begin{aligned}
 D_{\mathrm{KL}}(\rho_X\Vert\sigma_X)
 &=\int\phi(f\circ\widehat q)\,d\sigma_X\\
 &=\int\phi(f)\,d\sigma_M
 =D_{\mathrm{KL}}(\rho_M\Vert\sigma_M).
 \end{aligned}
\]
The Radon--Nikodym densities are finite almost everywhere; the
integrals may still be infinite. No integrability assumption or
reverse conditional kernel is needed.
\leanproof{MonadFormal.Information.klDiv_input_joint}
\end{proof}
